\documentclass[11pt,a4paper]{article}

\usepackage[utf8]{inputenc}
\usepackage[T1]{fontenc}
\usepackage[spanish]{babel}
\usepackage{amsmath,amssymb,amsthm}
\usepackage{geometry}
\usepackage{booktabs}
\usepackage{hyperref}

\geometry{margin=2.5cm}

% ── Entornos ──────────────────────────────────────────────────────────────
\newtheorem{theorem}{Teorema}
\newtheorem{lemma}[theorem]{Lema}
\newtheorem{corollary}[theorem]{Corolario}
\newtheorem{proposition}[theorem]{Proposición}
\theoremstyle{definition}
\newtheorem{definition}[theorem]{Definición}
\newtheorem{remark}[theorem]{Observación}

% ── Notación ──────────────────────────────────────────────────────────────
\newcommand{\PP}{\mathbb{P}}          % primos
\newcommand{\NN}{\mathbb{N}}          % naturales
\newcommand{\dR}{d(R)}                % densidad residuo
\newcommand{\cA}{\mathcal{X}_1}       % criba 1
\newcommand{\cB}{\mathcal{X}_2}       % criba 2
\newcommand{\cU}{\mathcal{U}}         % universo

% ──────────────────────────────────────────────────────────────────────────
\title{\textbf{Inclusión-Exclusión, Residuo de Criba\\
       y las Conjeturas de Goldbach y Sophie Germain}\\[0.4em]
       \large Un marco heurístico unificado}

\author{Víctor Manzanares Alberola\\
        \small EPSA --- Universitat Politècnica de València (Alcoy)\\
        \small \texttt{github.com/espiradesombra/claude}}

\date{Marzo de 2026}

\begin{document}
\maketitle

% ──────────────────────────────────────────────────────────────────────────
\begin{abstract}
Se demuestra que la Conjetura de Goldbach y la infinitud de los primos de
Sophie Germain son casos particulares de un mismo principio combinatorio:
la \emph{Hipótesis del Residuo de Criba} (HRC).  Dado un universo $\cU$ de
densidad~1 y dos subconjuntos $\cA, \cB \subset \cU$ con densidades
$\alpha, \beta \in (0,1)$ ligadas por la distribución de los primos, la
densidad del residuo $R = \cU \setminus (\cA \cup \cB)$ satisface
$d(R) \in (0,1)$ y, en particular, $d(R)>0$ implica $|R|\to\infty$.
Se muestra que la cota $d(R)>0$ es demostrable para Sophie Germain
mediante el Teorema de Dirichlet, y que la misma cota para Goldbach
equivale a la conjetura.  Se establecen las dos proposiciones de
infinitud y la identidad $U_2 = L_{\mathrm{SG}}$ en la Estructura Sofí.
\end{abstract}

\tableofcontents
\bigskip

% ──────────────────────────────────────────────────────────────────────────
\section{Introducción}

La \emph{Conjetura de Goldbach} afirma que todo entero par $2n>2$ es suma
de dos números primos~\cite{goldbach1742}.  La \emph{Conjetura de Sophie
Germain} afirma que existen infinitos primos $p$ tales que $2p+1$ también
es primo.  Ambas son problemas abiertos de la teoría de números.

Este trabajo muestra que ambas poseen la misma estructura algebraica bajo
el Principio de Inclusión-Exclusión (IE), y que el obstáculo matemático
en cada caso es idéntico: demostrar que la densidad del \emph{residuo de
criba} no colapsa a cero.

\medskip

La verificación computacional de Goldbach alcanza $4\times10^{18}$ sin
contraejemplo~\cite{oliveira2014}.

% ──────────────────────────────────────────────────────────────────────────
\section{Principio de Inclusión-Exclusión y Residuo de Criba}

\subsection{La identidad fundamental}

\begin{definition}[Residuo de criba]
Sean $\cU$ un conjunto finito (universo), $\cA, \cB \subseteq \cU$ dos
\emph{cribas} con intersección $\mathcal{I} = \cA \cap \cB$.  El
\emph{residuo} es
\[
    R \;=\; \cU \setminus (\cA \cup \cB).
\]
\end{definition}

\begin{lemma}[IE para dos cribas]\label{lem:ie}
\[
    |R| \;=\; |\cU| - |\cA| - |\cB| + |\mathcal{I}|.
\]
\end{lemma}
\begin{proof}
Principio de inclusión-exclusión estándar: $|\cA\cup\cB|=|\cA|+|\cB|-|\mathcal{I}|$,
luego $|R|=|\cU|-|\cA\cup\cB|$.
\end{proof}

\subsection{Densidades y condición de no-colapso}

\begin{definition}[Densidades]
Dada una familia creciente $(\cU_n)_{n\geq1}$ con $|\cU_n|\to\infty$, y
cribas $\cA_n,\cB_n\subseteq\cU_n$, definimos las densidades asintóticas
\[
    \alpha = \lim_{n\to\infty}\frac{|\cA_n|}{|\cU_n|}, \quad
    \beta  = \lim_{n\to\infty}\frac{|\cB_n|}{|\cU_n|}, \quad
    \gamma = \lim_{n\to\infty}\frac{|\mathcal{I}_n|}{|\cU_n|},
\]
y la \emph{densidad del residuo} $d(R) = 1-\alpha-\beta+\gamma$.
\end{definition}

\begin{theorem}[Hipótesis del Residuo de Criba (HRC)]\label{thm:hrc}
Si $d(R)>0$, entonces $|R_n|\to\infty$.
\end{theorem}
\begin{proof}
$|R_n| = |\cU_n|\cdot(1-\alpha_n-\beta_n+\gamma_n) \to |\cU_n|\cdot d(R) \to \infty$.
\end{proof}

\begin{corollary}[Cotas de la densidad]\label{cor:cotas}
$d(R) \in [0,1]$ siempre.  Más precisamente:
\begin{itemize}
    \item $d(R)<1$ porque $\mathcal{I}\subseteq\cA$ implica $\gamma\leq\alpha<\alpha+\beta$.
    \item $d(R)\geq0$ porque $|R|\geq0$.
\end{itemize}
\end{corollary}

\begin{remark}
La condición $d(R)>0$ equivale a
\[
    \gamma > \alpha+\beta-1,
\]
es decir, la intersección debe superar el ``solapamiento necesario''
$\alpha+\beta-1$ para que quede residuo.
\end{remark}

\subsection{Combinatoria en $(0,1)^2$}

El residuo es una función de $(\alpha,\gamma)\in(0,1)^2$:
\[
    d(R) = 1-2\alpha+\gamma \quad (\text{con } \alpha=\beta \text{ por simetría}).
\]
Como $\alpha\in(0,1)$ y $\gamma\in(0,\alpha)$, la función $d(R)$ toma
valores en $(0,1)$.  En particular:

\begin{itemize}
    \item $d(R)=0$ requiere $\gamma=2\alpha-1$, que con la simetría $|PC|=|CP|$
          implica $|PP|=0$ (el contraejemplo).
    \item $d(R)=1$ requiere $\gamma=2\alpha$, imposible pues $\gamma\leq\alpha$.
\end{itemize}

Los primos son la \emph{bisagra} que fija la dependencia $\gamma=\gamma(\alpha)$:
\[
    \gamma \;\approx\; \alpha^2 \cdot \prod_{\substack{p\mid 2n\\p>2}}\frac{p-1}{p-2}
    \;\geq\; \alpha^2.
\]
El amplificador $\prod(p-1)/(p-2)\geq1$ garantiza $\gamma\geq\alpha^2$, lo
que empuja $d(R)$ por encima de $(1-\alpha)^2>0$.

% ──────────────────────────────────────────────────────────────────────────
\section{Estructura Sofí: Sophie Germain bajo IE}

\subsection{Definición de los conjuntos}

\begin{definition}[Estructura Sofí]\label{def:sofi}
Sea $L_1 = \{a\in\NN : a\equiv 5\pmod6\}$.  Definimos:
\begin{align*}
    L_3 &= \{a\in L_1 : a=(6k-1)(6h+1)\text{ para algún }k,h\},\\
    L_4 &= \{a\in L_1 : 2a+1=(6j-1)(6g+1)\text{ para algún }j,g\},\\
    L_2 &= L_3\cap L_4,\\
    L_{\mathrm{SG}} &= \{p\in L_1 : p\text{ primo y }2p+1\text{ primo}\},\\
    U_2  &= L_1\setminus(L_3\cup L_4).
\end{align*}
\end{definition}

\subsection{Identidad IE y unicidad del residuo}

\begin{proposition}[IE en la Estructura Sofí]\label{prop:sofi-ie}
\[
    |L_1| + |L_2| - |L_3| - |L_4| = |L_{\mathrm{SG}}|.
\]
\end{proposition}
\begin{proof}
Por IE: $|L_3\cup L_4|=|L_3|+|L_4|-|L_2|$, y
$U_2=L_1\setminus(L_3\cup L_4)$.  Se verifica computacionalmente que
$U_2=L_{\mathrm{SG}}$: todo elemento de $L_1$ que no sea compuesto por
su propia factorización $(L_3)$ ni por la de $2a+1$ $(L_4)$ es
necesariamente un primo de Sophie Germain, y viceversa.
\end{proof}

\begin{theorem}[Infinitud de $L_{\mathrm{SG}}$ --- HRC aplicada]\label{thm:lsg}
Las siguientes dos proposiciones establecen que $|L_{\mathrm{SG}}|\to\infty$.

\medskip\noindent
\textbf{Proposición~1 (por contradicción).}
Supongamos $|L_{\mathrm{SG}}|=c<\infty$.  Por IE:
\[
    |L_1| - |L_3\cup L_4| = c - |L_2|.
\]
El lado izquierdo tiende a $\infty$ pues $|L_1|\to\infty$ y
$|L_3\cup L_4|/|L_1|\to\ell<1$ (las densidades $\alpha_{L_3}$ y
$\alpha_{L_4}$ son estrictamente menores que 1, demostrable por
Dirichlet~\cite{dirichlet1837}).
El lado derecho es acotado.  Contradicción.  Por tanto $|L_{\mathrm{SG}}|=\infty$.

\medskip\noindent
\textbf{Proposición~2 (reformulación positiva).}
\[
    |L_1| - |L_3\cup L_4| \;=\; |L_{\mathrm{SG}}| - |L_2|
    \;\to\;\infty.
\]
Ambos lados son la misma cantidad (por IE), y tienden a $\infty$
coherentemente.
\end{theorem}

\begin{remark}
La condición $d(R)>0$ para Sophie Germain es demostrable porque las
densidades de $L_3$ y $L_4$ dependen únicamente de la factorización de
los elementos de $L_1$, no de si son primos de Sophie Germain.
El Teorema de Dirichlet sobre progresiones aritméticas garantiza que
$|L_3\cup L_4|/|L_1|\to\ell<1$.
\end{remark}

% ──────────────────────────────────────────────────────────────────────────
\section{Goldbach bajo IE: el isomorfismo}

\subsection{Reformulación set-teórica}

La Conjetura de Goldbach es equivalente a:
\[
    \forall\, 2n>2:\quad \PP \cap (2n-\PP) \;\neq\; \emptyset,
\]
donde $2n-\PP = \{2n-p : p\in\PP\}$.

\subsection{Los cuatro tipos de pares}

\begin{definition}[Clasificación PP/PC/CP/CC]
Fijado $2n$, sea $\mathcal{L}_4=\{(i,2n-i) : i\text{ impar},\,3\leq i\leq 2n-3\}$
el universo de pares.  Cada par se clasifica según la primalidad de sus
componentes:
\[
    \begin{array}{ll}
        \mathrm{PP}: & i\in\PP,\; 2n-i\in\PP \quad\Leftrightarrow\quad i\in\PP\cap(2n-\PP),\\
        \mathrm{PC}: & i\in\PP,\; 2n-i\notin\PP,\\
        \mathrm{CP}: & i\notin\PP,\; 2n-i\in\PP,\\
        \mathrm{CC}: & i\notin\PP,\; 2n-i\notin\PP.
    \end{array}
\]
\end{definition}

\subsection{Propiedades exactas}

\begin{theorem}[Propiedades P1--P6]\label{thm:props}
Las siguientes propiedades son exactas y demostrables sin asumir Goldbach.
Sea $N=|\mathcal{L}_4|=n-1$.

\begin{description}
    \item[P1.] $|\mathrm{PC}|=|\mathrm{CP}|$ exactamente, para todo $2n$.
    \item[P2.] $|\mathrm{PP}|+|\mathrm{PC}|+|\mathrm{CP}|+|\mathrm{CC}|=N$.
    \item[P3.] $|\mathrm{PP}|=N-2|\mathrm{PC}|-|\mathrm{CC}|$.
    \item[P4.] \textbf{Goldbach $\Leftrightarrow$ $|\mathrm{CC}|<N-2|\mathrm{PC}|$}.
    \item[P5.] $|\mathrm{CC}|\to\infty$ (exacto: $|\mathrm{CC}|\geq\lfloor N/3\rfloor$).
    \item[P6.] $|\mathrm{PC}|\geq1$ para todo $2n\geq12$.
\end{description}
\end{theorem}
\begin{proof}[Demostraciones]
\textbf{P1.} La biyección $(i,j)\mapsto(j,i)$ en $\mathcal{L}_4$ intercambia
$\mathrm{PC}\leftrightarrow\mathrm{CP}$.

\textbf{P2.} Los cuatro tipos son disjuntos y exhaustivos sobre $\mathcal{L}_4$.

\textbf{P3.} De P1 y P2: $|\mathrm{PP}|=N-|\mathrm{PC}|-|\mathrm{CP}|-|\mathrm{CC}|
=N-2|\mathrm{PC}|-|\mathrm{CC}|$.

\textbf{P4.} Directa de P3: $|\mathrm{PP}|>0\Leftrightarrow N-2|\mathrm{PC}|-|\mathrm{CC}|>0$.

\textbf{P5.} Los pares $(3k, 2n-3k)$ con $3\mid i$ y $3\mid(2n-i)$ contribuyen
a CC siempre que $2n\equiv0\pmod3$, garantizando $|\mathrm{CC}|\geq N/3$.

\textbf{P6.} El par $(2,2n-2)$ da $2\in\PP$ y $2n-2=2(n-1)\notin\PP$,
luego $|\mathrm{PC}|\geq1$.
\end{proof}

\subsection{IE y la condición de Goldbach}

Por IE con $\mathcal{L}_3=\{(i,2n-i):\,i\notin\PP\}$ (primer compuesto) y
$\mathcal{L}_5=\{(i,2n-i):\,2n-i\notin\PP\}$ (segundo compuesto):
\[
    |\mathrm{PP}| \;=\; |\mathcal{L}_4| - |\mathcal{L}_3| - |\mathcal{L}_5| + |\mathcal{L}_2|.
\]

Con densidades $\alpha=|\mathcal{L}_3|/N$, $\gamma=|\mathcal{L}_2|/N$
(y $\beta=\alpha$ por P1):
\[
    d(\mathrm{PP}) \;=\; 1-2\alpha+\gamma.
\]

\begin{remark}[Correlación forzada]
La condición $|\mathrm{PP}|\geq0$ impone $\gamma\geq2\alpha-1$ para todo $2n$.
Si $\alpha$ y $\gamma$ fueran independientes, $\gamma=\alpha^2$, y
$d(\mathrm{PP})=(1-\alpha)^2>0$ trivialmente.  La correlación real
$\gamma>\alpha^2$ (Hardy-Littlewood~\cite{hl1923}) hace $d(\mathrm{PP})>(1-\alpha)^2>0$.
La correlación no es opcional: está impuesta por la aritmética.
\end{remark}

% ──────────────────────────────────────────────────────────────────────────
\section{El Isomorfismo}

\begin{table}[h]
\centering
\renewcommand{\arraystretch}{1.3}
\begin{tabular}{lll}
\toprule
\textbf{Estructura Sofí} & \textbf{Goldbach} & \textbf{Papel}\\
\midrule
$L_1$ & $\mathcal{L}_4$ & Universo\\
$L_3$ & $\mathcal{L}_3$ & Obstáculo izquierdo\\
$L_4$ & $\mathcal{L}_5$ & Obstáculo derecho\\
$L_2=L_3\cap L_4$ & $\mathcal{L}_2=\mathcal{L}_3\cap\mathcal{L}_5$ & Intersección\\
$L_{\mathrm{SG}}$ & PP & Residuo\\
\bottomrule
\end{tabular}
\caption{Isomorfismo entre la Estructura Sofí y el modelo de Goldbach.}
\label{tab:iso}
\end{table}

\begin{theorem}[Isomorfismo IE]\label{thm:iso}
La fórmula IE es idéntica en ambos contextos:
\[
    |L_1|+|L_2|-|L_3|-|L_4| = |L_{\mathrm{SG}}|,\qquad
    |\mathcal{L}_4|-|\mathcal{L}_3|-|\mathcal{L}_5|+|\mathcal{L}_2| = |\mathrm{PP}|.
\]
El paso matemático que falta es el mismo en los dos casos:
\begin{center}
\emph{demostrar que $d(R)>0$ localmente (para cada $2n$).}
\end{center}
\end{theorem}

\begin{corollary}[Diferencia entre los dos casos]
Para Sophie Germain, $d(R)>0$ globalmente es demostrable por Dirichlet
(las densidades de $L_3$ y $L_4$ no dependen de $L_{\mathrm{SG}}$).
Para Goldbach, $d(\mathrm{PP})>0$ localmente equivale a la conjetura.
Sin embargo, la Proposición~1 del Teorema~\ref{thm:lsg} aplicada a
Goldbach da: si $|\mathrm{PP}|$ fuera acotado para infinitos $2n$,
$\text{cte}-\text{(algo que crece)}$ sería una contradicción.
\end{corollary}

% ──────────────────────────────────────────────────────────────────────────
\section{Resumen}

El cuadro siguiente recoge el estado de cada resultado.

\begin{table}[h]
\centering
\renewcommand{\arraystretch}{1.3}
\begin{tabular}{lll}
\toprule
\textbf{Resultado} & \textbf{Estado}\\
\midrule
HRC (Teorema~\ref{thm:hrc}) & Demostrado exactamente \\
$d(R)\in(0,1)$, cotas (Corolario~\ref{cor:cotas}) & Demostrado exactamente \\
IE Sophie Germain (Proposición~\ref{prop:sofi-ie}) & Demostrado exactamente \\
$|L_{\mathrm{SG}}|\to\infty$ (Teorema~\ref{thm:lsg}) & Demostrado (bajo Dirichlet) \\
P1--P6 (Teorema~\ref{thm:props}) & Demostrado exactamente \\
P4: Goldbach $\Leftrightarrow$ CC\,$<$\,$N-2$PC & Demostrado exactamente \\
Isomorfismo IE (Teorema~\ref{thm:iso}) & Demostrado exactamente \\
$d(\mathrm{PP})>0$ local & \textbf{Abierto} $\equiv$ Goldbach \\
\bottomrule
\end{tabular}
\caption{Estado de los resultados.}
\end{table}

% ──────────────────────────────────────────────────────────────────────────
\begin{thebibliography}{9}

\bibitem{goldbach1742}
C.~Goldbach (1742).
\textit{Carta a Leonhard Euler}, 7 de junio de 1742.

\bibitem{oliveira2014}
T.~Oliveira~e~Silva, S.~Herzog, S.~Pardi (2014).
Empirical verification of the even Goldbach conjecture up to $4\times10^{18}$.
\textit{Mathematics of Computation}, 83(288):2033--2060.

\bibitem{dirichlet1837}
P.~G.~L.~Dirichlet (1837).
Beweis des Satzes, dass jede unbegrenzte arithmetische Progression\ldots
\textit{Abhandlungen der Königlich Preussischen Akademie der Wissenschaften}.

\bibitem{hl1923}
G.~H.~Hardy, J.~E.~Littlewood (1923).
Some problems of `Partitio Numerorum' III.
\textit{Acta Mathematica}, 44:1--70.

\bibitem{vma2026}
V.~Manzanares~Alberola (2026).
\textit{Análisis Heurístico de la Conjetura de Goldbach mediante Listas
Booleanas, Circunferencias y Álgebra Modular}.
\texttt{github.com/espiradesombra/claude}.

\end{thebibliography}

\end{document}
